(1) 理想气体质量密度为 \(\rho=Mp/(RT)\)，静力平衡是 \(dp/dz=-\rho g\)，因此
\[
\frac{dp}{dz}=-\frac{Mgp}{RT}.
\]
该式对局部温度随高度变化完全一般，不要求大气等温或绝热。

(2) 理想气体每摩尔熵微分为 \(ds=C_p\,dT/T-R\,dp/p\)。随小气团可逆绝热升降令 ds=0，于是
\[
\frac{dT}{dz}=\frac{R}{C_p}\frac Tp\frac{dp}{dz}
=-\frac{Mg}{C_p},
\quad \Gamma_d=\frac{Mg}{C_p}.
\]
定义位温 \(\vartheta=T(p_0/p)^{R/C_p}\)，同一空气团绝热位移时 \(\vartheta\) 保持不变；位温升高意味着在同一参考压强下的温度升高，也对应熵升高。

(3) 令背景温度直减率为 Γ，空气团相对初始平衡位置向上位移 ξ>0。空气团绝热冷却 \(\delta T_{\rm parcel}=-\Gamma_d\xi\)，同一高度环境变化 \(\delta T_{\rm env}=-\Gamma\xi\)。两者压强相同，因此密度差在一阶近似下
\(\delta\rho/\rho\simeq-(T_{\rm parcel}-T_{\rm env})/T
=(\Gamma_d-\Gamma)\xi/T\)。
向上浮力加速度为 \(-g\delta\rho/\rho\)，故
\[
\ddot\xi+N^2\xi=0,\qquad N^2=\frac gT(\Gamma_d-\Gamma).
\]

(4) 若 Γ<Γd，N²>0，小气团在原位附近作 Brunt–Väisälä 稳定浮力振荡，周期为 \(2\pi/N\)。若 Γ>Γd，N²<0，扰动出现指数增长，对流不稳定；Γ=Γd 时线性回复力为零，处于中性绝热层结。以上结果依赖于压强快速平衡、干绝热和局部小位移的假设，水汽凝结会改变有效直减率。
